\section{Introduction}
\label{m:sec:intro}

The processes studied in this thesis describe populations small enough that their
fate is decided by chance rather than by their mean behaviour --- a founding cohort
of a few individuals, or the contents of a single compartment that may later
rupture. At those counts a deterministic rate equation answers the wrong question.
It returns a single trajectory, growing or decaying, when what one wants to know is
whether the population establishes itself at all, and how large it is if it does.
Both are questions about a distribution, and neither is visible in a mean.

This chapter assembles the probabilistic and analytic toolkit that the rest of the
thesis uses. Two threads run through it, and they are worth separating at the
outset, because almost every later section belongs to one or the other.

The first is \emph{conditioning on survival}. A subcritical branching process dies
out with probability one, so its unconditional mean decays geometrically and
records little beyond the certainty of extinction. Restrict attention to those
realisations that happen still to be alive at generation $n$, however, and the
picture changes entirely: the conditional population settles to a fixed
distribution, and the mean of that distribution is a constant. For the binary
Galton--Watson process with replication probability $p<\tfrac12$ the constant is
$1/A(p)$, where
\begin{equation*}
A(p) = \lim_{n\to\infty}\frac{S_n}{(2p)^n}
\end{equation*}
and $S_n$ is the probability of survival to generation $n$. This chapter defines
$A(p)$, establishes that the limit exists, and shows how the whole of the limiting
conditional law is expressed through it. It does not analyse the constant itself:
the infinite product, the exact series and its bounds, the near-critical
asymptotic, the search for a closed form, and the dynamical obstruction that
explains why that search failed are the subject of \ChA.

The second thread is the \emph{near-critical regime}. As $p$ rises towards
$\tfrac12$ the process takes ever longer to die, the quasi-stationary mean $1/A(p)$
diverges, and any numerical scheme that works by iterating to stationarity becomes
impractical. Just above $\tfrac12$ the size of the founding cohort begins to weigh
heavily on whether the process survives at all. The behaviour on either side of
$p=\tfrac12$ is what makes small founding populations interesting, and it is also
what makes them awkward to compute with. The amplitude theory that eventually
tames it is again the business of \ChA; what belongs here is the phenomenology and
the continuous-time comparison, where a closed form \emph{is} available.

\Cref{m:sec:prelim} assembles the standard apparatus: exponential clocks and the
races between them, continuous-time Markov chains and their generators,
probability generating functions, and the linear birth--death process, whose
generating-function partial differential equation is the template for everything in
\cref{m:sec:moc}. \Cref{m:sec:gw} sets out the Galton--Watson process and
establishes the two facts on which the rest depends --- that extinction is certain
at and below criticality, and that at criticality the expected lifetime is
nevertheless infinite. \Cref{m:sec:qsd} derives the limiting conditional mean and
variance, and so introduces $A(p)$. \Cref{m:sec:small} asks what changes when the
process starts from a cohort of $k$ individuals rather than one, treats a discrete
birth--death process subject to catastrophe, obtains the corresponding constant
$\Ac(p)$ for the continuous-time birth--death process in closed form, and defines
the continuous-time birth--death--catastrophe process used in later chapters.
\Cref{m:sec:moc} then changes tack and develops the method of characteristics for
the generating-function equations of absorption models.

Two remarks on how the chapter is written. First, not all of it is background. The
closed form obtained in \cref{m:sec:absbirthdeath} for the
absorption--birth--death generating function is a result of this thesis, as is the
integral identity of \cref{m:app:hyp} on which an alternative derivation rests;
everything else in the chapter is standard, and is presented as such. Second, the
argument is in places heuristic --- continuum approximations to difference
equations, mean-field substitutions, asymptotic matching. Every such step is marked
where it occurs, and where a claim can be checked numerically it has been.
